% Seção 4 — Objetivos e hipóteses. Reescrita em 2026-11-15.

\chapter{Objetivos e hipóteses}
\label{cap:objetivos}

\textbf{Objetivo geral.} Estimar o efeito causal da proibição estadual da
pulverização aérea de agrotóxicos (Lei nº~16.820/2019) sobre a saúde ao nascer
nos municípios cearenses onde a prática existia.

\textbf{Objetivos específicos.}
\begin{alineas}
  \item medir onde e quando havia aplicação aérea antes da proibição, com
    registros administrativos federais e estaduais;
  \item verificar, só com dados anteriores à lei, se a exposição à estação de
    pulverização deixava marca sazonal no peso ao nascer;
  \item estimar o efeito da proibição pelo desenho de calendário gestacional,
    com os municípios de banana familiar como grupo de comparação;
  \item confrontar essa estimativa com a forma reduzida estadual, o controle
    sintético e a comparação de fronteira;
  \item testar a seleção para o nascimento vivo e o canal de substituição da
    aplicação aérea pela terrestre.
\end{alineas}

\textbf{Hipóteses.} A hipótese principal ($H_1$) é que, após a lei, o peso ao
nascer das coortes gestadas na estação de pulverização aumentou, em relação às
demais coortes, nos municípios onde havia aplicação aérea, comparados aos de
banana familiar. A segunda hipótese confirmatória ($H_2$) é que a taxa de
óbito fetal dessas coortes caiu. Como hipóteses de mecanismo, o efeito deve
ser maior nos anos chuvosos ($H_3$) e ausente quando só o terceiro trimestre
coincide com a estação e nos municípios de banana familiar comparados aos sem
banana ($H_4$). Como hipótese de substituição ($H_5$), as notificações de
intoxicação ocupacional por agrotóxico agrícola podem aumentar, se a aplicação
migrou do avião para o trator e o pulverizador costal.
